\documentclass[12pt]{article}
\usepackage{amsmath, amssymb, amsthm}
\usepackage{hyperref}

\newtheorem{theorem}{Theorem}
\newtheorem{lemma}[theorem]{Lemma}
\newtheorem{proposition}[theorem]{Proposition}
\newtheorem{corollary}[theorem]{Corollary}
\newtheorem{definition}[theorem]{Definition}
\newtheorem{remark}[theorem]{Remark}

\title{Problem 10: Preconditioned CG for RKHS-Constrained CP Decomposition\\
\large A Complete Explanation}
\author{}
\date{April 2026}

\begin{document}
\maketitle

\section{Problem Statement}

We have a $d$-way tensor $\mathcal{T} \in \mathbb{R}^{n_1 \times \cdots \times n_d}$ with missing entries, and we seek a rank-$r$ CP decomposition where mode $k$ is constrained to a Reproducing Kernel Hilbert Space (RKHS). The mode-$k$ subproblem (with all other factor matrices fixed) requires solving a linear system of size $nr \times nr$:
\begin{equation}\label{eq:system}
\underbrace{\left[(Z \otimes K)^T S\, S^T (Z \otimes K) + \lambda(I_r \otimes K)\right]}_{\mathcal{A}} \mathrm{vec}(W) = (I_r \otimes K)\,\mathrm{vec}(B).
\end{equation}

Here:
\begin{itemize}
\item $K \in \mathbb{R}^{n \times n}$: PSD RKHS kernel matrix (PSD, symmetric).
\item $Z = A_d \odot \cdots \odot A_{k+1} \odot A_{k-1} \odot \cdots \odot A_1 \in \mathbb{R}^{M \times r}$: Khatri-Rao product.
\item $S \in \mathbb{R}^{N \times q}$: selection matrix (a subset of $N$ identity columns), $q \ll N$.
\item $B = TZ \in \mathbb{R}^{n \times r}$: the MTTKRP.
\item $\lambda > 0$: regularization parameter.
\item $n, r \ll M \ll N$.
\end{itemize}

We want to solve \eqref{eq:system} efficiently using preconditioned conjugate gradient (PCG), avoiding any $O(N)$ computation.

\section{Setup and Notation}

\subsection{Matrix decomposition of the system}

Let $\mathcal{A} = (Z \otimes K)^T S S^T (Z \otimes K) + \lambda (I_r \otimes K)$.

We first simplify. Let $\hat{S} = S^T \in \mathbb{R}^{q \times N}$ be the row-selection matrix. Then $SS^T$ is the projection onto the observed entries. The matrix $Z \otimes K \in \mathbb{R}^{nM \times nr}$... wait, let me recheck dimensions.

\textbf{Dimension check:}
\begin{itemize}
\item $Z \in \mathbb{R}^{M \times r}$, so $Z \otimes K \in \mathbb{R}^{(nM) \times (nr)}$? No: $(Z \otimes K) \in \mathbb{R}^{nM \times nr}$... let me be more careful.

Actually: $(Z \otimes K)$ should be of size $N \times nr$ where $N = nM$ (since the mode-$k$ unfolding $T$ is $n \times M$ and $\mathrm{vec}(T)$ is of length $nM = N$). But the problem says $S \in \mathbb{R}^{N \times q}$, confirming $N = nM$.

However, the Khatri-Rao product $Z \in \mathbb{R}^{M \times r}$ and $I_n \otimes Z$ would be $nM \times nr$... let me re-read.

The system matrix is $(Z \otimes K)^T S S^T (Z \otimes K)$. With $(Z \otimes K) \in \mathbb{R}^{N \times nr}$ (this requires $N = nr$?).

Actually, from the problem: $N = \prod_i n_i$ is the total tensor size, $M = \prod_{i \neq k} n_i$, $n = n_k$. The mode-$k$ unfolding $T \in \mathbb{R}^{n \times M}$ has $\mathrm{vec}(T) \in \mathbb{R}^{nM} = \mathbb{R}^N$.

The CP model says $\mathrm{vec}(T) \approx (Z \otimes I_n) \mathrm{vec}(A_k)$ (in the non-RKHS case). With RKHS: $A_k = KW$, so $\mathrm{vec}(A_k) = \mathrm{vec}(KW) = (I_r \otimes K)\mathrm{vec}(W)$.

So the model is $\mathrm{vec}(T) \approx (Z \otimes I_n)(I_r \otimes K)\mathrm{vec}(W) = (Z \otimes K)\mathrm{vec}(W)$.

Here $(Z \otimes K) \in \mathbb{R}^{NM... no.}$

Let me redo: $Z \in \mathbb{R}^{M \times r}$, $K \in \mathbb{R}^{n \times n}$. The Kronecker product $Z \otimes K \in \mathbb{R}^{(Mr) \times (nM)}$? No.

$Z \otimes K = $ the matrix with $(Z)_{ij}$ replaced by $(Z)_{ij} K$, giving size $(M \cdot n) \times (r \cdot n)$... that's $Mn \times nr = N \times nr$. Wait but $S \in \mathbb{R}^{N \times q}$, so $S^T(Z \otimes K) \in \mathbb{R}^{q \times nr}$, and $(Z\otimes K)^T S S^T (Z\otimes K) \in \mathbb{R}^{nr \times nr}$.

Yes: the system \eqref{eq:system} is of size $nr \times nr$ as stated.

Also, $\lambda(I_r \otimes K) \in \mathbb{R}^{nr \times nr}$ (correct, since $I_r \in \mathbb{R}^{r\times r}$, $K \in \mathbb{R}^{n\times n}$, so $I_r \otimes K \in \mathbb{R}^{nr \times nr}$).

So the system is $\mathcal{A} \cdot \mathrm{vec}(W) = (I_r \otimes K)\mathrm{vec}(B)$ where $\mathcal{A} \in \mathbb{R}^{nr \times nr}$.
\end{itemize}

\subsection{Observed entries structure}

The selection matrix $S \in \mathbb{R}^{N \times q}$ selects the $q$ observed tensor entries. The matrix $SS^T \in \mathbb{R}^{N \times N}$ is diagonal with 1's at observed positions and 0's elsewhere.

We write $SS^T = \Omega$ (the observation operator / projection onto observed entries).

\section{Preconditioned Conjugate Gradient Method}

\subsection{Why use PCG?}

Direct solution of \eqref{eq:system} using Cholesky decomposition costs $O((nr)^3) = O(n^3 r^3)$ (as noted in the problem). For large $n$ and $r$, this is prohibitive.

\textbf{The CG approach:} The conjugate gradient method for a PD system $\mathcal{A}x = b$ requires only matrix-vector products with $\mathcal{A}$. We show these can be done in $O(qnr)$ time (linear in the number of observations times the CP rank times the number of RKHS points).

\subsection{The preconditioner}

\begin{definition}[Preconditioner]\label{def:preconditioner}
A preconditioner $P$ for the system $\mathcal{A}x = b$ is a symmetric positive definite matrix such that:
\begin{enumerate}
\item $P \approx \mathcal{A}$ (in some sense), so the condition number $\kappa(P^{-1}\mathcal{A})$ is small.
\item $Py = v$ can be solved cheaply (ideally in $O(nr)$ or $O(nr \log(nr))$ time).
\end{enumerate}
\end{definition}

\textbf{Proposed preconditioner:}

We take:
\begin{equation}\label{eq:precond}
P = \lambda (I_r \otimes K) + \mathrm{diag}(\mathcal{A} - \lambda(I_r \otimes K)),
\end{equation}
or more practically:
\begin{equation}\label{eq:precond2}
P = \lambda (I_r \otimes K).
\end{equation}

The preconditioner $P = \lambda(I_r \otimes K)$ is the RKHS regularization term alone. It is PD since $K \succ 0$ (by assumption) and $\lambda > 0$.

\textbf{Solving $Py = v$:} This requires $(I_r \otimes K)^{-1} v = (I_r \otimes K^{-1})v$, i.e., applying $K^{-1}$ to each of the $r$ blocks of the $nr$-vector $v$. Cost: $O(n^2 r)$ (if $K^{-1}$ is precomputed, or $O(n^3 + n^2 r)$ total if Cholesky of $K$ is precomputed).

\textbf{Better preconditioner:}

A better preconditioner exploits the Kronecker structure:
\begin{equation}\label{eq:precond-kron}
P = (Z^T Z + \lambda I_r) \otimes K,
\end{equation}
which approximates $\mathcal{A}$ by replacing $\Omega = SS^T$ with the identity. When all entries are observed ($q = N$, $S = I_N$):
\[
(Z \otimes K)^T (Z \otimes K) + \lambda(I_r \otimes K) = (Z^T Z \otimes K^T K) + \lambda(I_r \otimes K) = (Z^T Z \otimes K^2) + \lambda(I_r \otimes K).
\]
Hmm, this doesn't factor as nicely. Let me recompute.

$(Z \otimes K)^T = Z^T \otimes K^T = Z^T \otimes K$ (since $K$ is symmetric). So:
\[
(Z\otimes K)^T(Z\otimes K) = (Z^T \otimes K)(Z \otimes K) = (Z^T Z) \otimes K^2.
\]

And $\mathcal{A}_{\Omega=I} = (Z^T Z) \otimes K^2 + \lambda (I_r \otimes K)$.

A Kronecker-structured preconditioner for this is $P = (Z^T Z) \otimes K^2 + \lambda(I_r \otimes K)$... but applying $P^{-1}$ requires simultaneous diagonalization of $Z^T Z$ and $K^2$, which is possible since both are symmetric:

Let $K = V \Sigma V^T$ (eigendecomposition, $V$ orthogonal, $\Sigma = \mathrm{diag}(\sigma_1,\ldots,\sigma_n)$) and $Z^T Z = U \Lambda U^T$ (eigendecomposition). Then:
\[
P = (U \otimes V)(\Lambda \otimes \Sigma^2 + \lambda I_{nr})(U \otimes V)^T.
\]
Applying $P^{-1}$:
\begin{itemize}
\item Apply $(U^T \otimes V^T)$: cost $O(n^2 r + nr^2)$.
\item Apply diagonal matrix $(\Lambda \otimes \Sigma^2 + \lambda I)^{-1}$: cost $O(nr)$.
\item Apply $(U \otimes V)$: cost $O(n^2 r + nr^2)$.
\end{itemize}
Total cost of applying $P^{-1}$: $O(n^2 r + nr^2)$ (after precomputing eigendecompositions of $K$ and $Z^T Z$).

\section{Matrix-Vector Products with $\mathcal{A}$}

The key computational challenge is applying $\mathcal{A}$ to a vector $v \in \mathbb{R}^{nr}$:
\[
\mathcal{A}v = (Z \otimes K)^T S S^T (Z \otimes K) v + \lambda(I_r \otimes K)v.
\]

\subsection{Step-by-step computation}

Let $v \in \mathbb{R}^{nr}$, viewed as $\mathrm{vec}(V)$ for $V \in \mathbb{R}^{n \times r}$.

\textbf{Term 1: $\lambda(I_r \otimes K)v$.}

$(I_r \otimes K)v = \mathrm{vec}(KV)$. This requires $r$ applications of $K \in \mathbb{R}^{n\times n}$ to $n$-vectors.
Cost: $O(n^2 r)$.

\textbf{Term 2: $(Z \otimes K)^T S S^T (Z \otimes K) v$.}

We compute from inside out:

\textbf{Step 2a: Compute $w_1 = (Z \otimes K)v$.}

$(Z \otimes K)v = \mathrm{vec}(KVZ^T)$.

To see this: $(Z \otimes K)\mathrm{vec}(V) = \mathrm{vec}(KVZ^T)$ by the identity $(A \otimes B)\mathrm{vec}(X) = \mathrm{vec}(BXA^T)$.

Cost: Apply $K$ to $r$ columns: $O(n^2 r)$. Then multiply by $Z^T$ (a $r \times M$ matrix applied to $n \times r$ result... wait.

Actually: $KV \in \mathbb{R}^{n \times r}$, and $KVZ^T \in \mathbb{R}^{n \times M}$. So $w_1 = \mathrm{vec}(KVZ^T) \in \mathbb{R}^{nM} = \mathbb{R}^N$.

Cost of forming $KV$: $O(n^2 r)$ (since $K$ is $n\times n$ and $V$ is $n \times r$).
Cost of forming $KVZ^T$: $O(nrM)$ (since $KV$ is $n\times r$ and $Z$ is $M\times r$, computing $(KV)Z^T = $ ($n\times r$) times ($r\times M$) matrix).

But $M = N/n$ is large ($M \ll N$ but $M$ can still be $\gg n, r$)! This gives $O(nrM) = O(rN/n \cdot n) = O(rN)$ which is $O(N)$ and violates our constraint.

\textbf{We do NOT explicitly form $w_1 = (Z\otimes K)v$.} Instead, we use the observation structure.

\textbf{Step 2b: Compute $w_2 = S^T w_1$.}

$S^T \mathrm{vec}(KVZ^T)$ selects $q$ entries from the $N$-vector $\mathrm{vec}(KVZ^T)$.

The $(i,j)$-entry of $KVZ^T$ (where $(i,j) \in [n] \times [M]$ corresponds to tensor index $(i, j_1, \ldots, j_{d-1})$... well, $j$ is the combined index over all modes except $k$) is:
\[
(KVZ^T)_{ij} = \sum_{\ell=1}^r (KV)_{i\ell} Z_{j\ell}.
\]

The selection matrix $S$ selects $q$ entries of $\mathrm{vec}(KVZ^T)$, i.e., $q$ pairs $(i_s, j_s)$ for $s = 1, \ldots, q$ (the observed entries in mode-$k$ unfolding coordinates).

So $w_2 \in \mathbb{R}^q$ has $s$-th entry:
\[
(w_2)_s = (KVZ^T)_{i_s j_s} = \sum_\ell (KV)_{i_s \ell} Z_{j_s \ell} = \langle (KV)_{i_s, :}, Z_{j_s, :}\rangle,
\]
i.e., the dot product of the $i_s$-th row of $KV$ with the $j_s$-th row of $Z$.

\textbf{Computing $w_2$ without forming $KVZ^T$:}

For each observed entry $s = 1, \ldots, q$:
\[
(w_2)_s = \sum_{\ell=1}^r (KV)_{i_s \ell} Z_{j_s \ell}.
\]

We first compute $KV \in \mathbb{R}^{n \times r}$ (cost: $O(n^2 r)$). Then for each observed entry $s$:
\[
(w_2)_s = \sum_\ell (KV)_{i_s \ell} Z_{j_s \ell}.
\]
Cost for all $q$ entries: $O(qr)$ (each entry requires $r$ multiplications).

Total cost of Step 2b: $O(n^2 r + qr)$.

\textbf{Step 2c: Compute $w_3 = S w_2$.}

$S w_2 \in \mathbb{R}^N$ is the vector that places $(w_2)_s$ at position corresponding to observed entry $s$ and $0$ elsewhere.

Conceptually, $w_3$ is the $N$-vector equal to $w_2$ on the $q$ observed coordinates and $0$ elsewhere.

We do not need to explicitly store $w_3$ as an $N$-vector. We only need $S^T (Z\otimes K) \mathcal{A}^{-1}\cdots$, which we compute directly.

\textbf{Step 2d: Compute $w_4 = (Z \otimes K)^T w_3$.}

$(Z\otimes K)^T w_3 = (Z^T \otimes K) w_3$ where $w_3 = S w_2$ is the sparse vector.

$(Z^T \otimes K) S w_2 \in \mathbb{R}^{nr}$.

Using the identity $(Z^T \otimes K)\mathrm{vec}(X) = \mathrm{vec}(KXZ)$ for $X \in \mathbb{R}^{n \times M}$:

$w_3 = Sw_2 = \mathrm{vec}(X_{\mathrm{obs}})$ where $X_{\mathrm{obs}} \in \mathbb{R}^{n \times M}$ is the sparse matrix with $(X_{\mathrm{obs}})_{i_s j_s} = (w_2)_s$ for each observed entry $s$.

Then $(Z^T \otimes K)\mathrm{vec}(X_{\mathrm{obs}}) = \mathrm{vec}(K X_{\mathrm{obs}} Z)$.

Computing $X_{\mathrm{obs}} Z \in \mathbb{R}^{n \times r}$: since $X_{\mathrm{obs}}$ has only $q$ nonzero entries:
\[
(X_{\mathrm{obs}} Z)_{i, \ell} = \sum_{j=1}^M (X_{\mathrm{obs}})_{ij} Z_{j\ell} = \sum_{s: i_s = i} (w_2)_s Z_{j_s \ell}.
\]

Cost: $O(qr)$ (for each of the $q$ observed entries, add one rank-1 contribution to the $n \times r$ matrix).

Then $K(X_{\mathrm{obs}} Z) \in \mathbb{R}^{n \times r}$: cost $O(n^2 r)$.

So $w_4 = \mathrm{vec}(K X_{\mathrm{obs}} Z) \in \mathbb{R}^{nr}$.

\textbf{Summary of matrix-vector product:}
\[
\mathcal{A}v = w_4 + \lambda(I_r \otimes K)v,
\]
with:
\begin{enumerate}
\item Compute $KV \in \mathbb{R}^{n \times r}$: $O(n^2 r)$.
\item Compute $(w_2)_s = \langle (KV)_{i_s,:}, Z_{j_s,:}\rangle$ for each observation $s$: $O(qr)$.
\item Form $X_{\mathrm{obs}} Z \in \mathbb{R}^{n\times r}$ by accumulating: $O(qr)$.
\item Compute $K(X_{\mathrm{obs}}Z) \in \mathbb{R}^{n\times r}$: $O(n^2 r)$.
\item Compute $K V$ (already done in step 1): $O(0)$.
\item Sum: $O(1)$.
\end{enumerate}

Total cost per matrix-vector product: $O(n^2 r + qr)$.

\section{Preconditioned CG Algorithm}

The full PCG algorithm for solving \eqref{eq:system}:

\begin{enumerate}
\item \textbf{Precomputation (done once):}
   \begin{itemize}
   \item Compute $Z^T Z \in \mathbb{R}^{r\times r}$: cost $O(Mr^2)$ (or use the structure of $Z$).
   \item Compute eigendecomposition $K = V_K \Sigma_K V_K^T$: cost $O(n^3)$.
   \item Compute eigendecomposition $Z^T Z = U_{Z} \Lambda_Z U_Z^T$: cost $O(r^3)$.
   \item Compute MTTKRP $B = TZ$: cost $O(qr)$ (only $q$ observed entries contribute).
   \item Compute right-hand side $b = (I_r \otimes K)\mathrm{vec}(B) = \mathrm{vec}(KB)$: cost $O(n^2 r)$.
   \end{itemize}

\item \textbf{PCG iteration:}
   \begin{itemize}
   \item Initialize $x_0 = 0$, $r_0 = b$, $p_0 = P^{-1}r_0$.
   \item At each iteration $t$:
      \begin{itemize}
      \item Compute $\mathcal{A} p_t$: cost $O(n^2 r + qr)$ (as above).
      \item Update $x_{t+1}$, $r_{t+1}$, $p_{t+1}$ by standard CG formulas.
      \item Apply $P^{-1}$ to $r_{t+1}$: cost $O(n^2 r + nr^2)$ (eigendecomposition approach).
      \end{itemize}
   \item Stop when $\|r_t\| / \|r_0\| < \varepsilon_{\mathrm{tol}}$.
   \end{itemize}
\end{enumerate}

\subsection{Why PCG converges rapidly}

\begin{theorem}[CG convergence bound]\label{thm:cg-conv}
The conjugate gradient method for $\mathcal{A}x = b$ (with $\mathcal{A}$ PD) satisfies:
\[
\|\mathcal{A}^{1/2}(x_t - x^*)\|_2 \leq 2\left(\frac{\sqrt{\kappa} - 1}{\sqrt{\kappa} + 1}\right)^t \|\mathcal{A}^{1/2}(x_0 - x^*)\|_2,
\]
where $\kappa = \kappa(\mathcal{A}) = \lambda_{\max}(\mathcal{A})/\lambda_{\min}(\mathcal{A})$ is the condition number.
\end{theorem}

\begin{proof}
This is the standard CG convergence theorem \cite[Theorem 6.29]{Trefethen1997}. The hypotheses are: $\mathcal{A}$ symmetric positive definite (verified: $\mathcal{A} = \mathcal{A}^T$ by construction, and $\mathcal{A} \succ 0$ because $\lambda(I_r \otimes K) \succ 0$ and the first term is PSD). The conclusion is the stated bound.
\end{proof}

For PCG with preconditioner $P$, the convergence rate is determined by $\kappa(P^{-1/2}\mathcal{A}P^{-1/2})$.

\begin{theorem}[Condition number analysis]\label{thm:cond}
With preconditioner $P = \lambda(I_r \otimes K)$:
\[
\kappa(P^{-1}\mathcal{A}) = 1 + \frac{1}{\lambda} \kappa\left((I_r \otimes K)^{-1} (Z\otimes K)^T SS^T (Z\otimes K)\right).
\]

More precisely:
\[
\lambda_{\max}(P^{-1}\mathcal{A}) = 1 + \frac{\sigma_{\max}((Z \otimes I_n)^T \Omega (Z \otimes I_n)) \sigma_{\max}(K)}{\lambda},
\]
where $\sigma_{\max}$ denotes the maximum eigenvalue and $\Omega = SS^T$ is the observation mask.
\end{theorem}

\begin{proof}
We compute $(P^{-1}\mathcal{A})v$ for an eigenvector $v$:
\[
P^{-1}\mathcal{A} = \frac{1}{\lambda}(I_r \otimes K^{-1})\left[(Z\otimes K)^T SS^T (Z\otimes K) + \lambda(I_r \otimes K)\right].
\]
\[
= \frac{1}{\lambda}(I_r \otimes K^{-1})(Z\otimes K)^T SS^T (Z\otimes K) + I_{nr}.
\]

The spectrum of $P^{-1}\mathcal{A}$ consists of $1$ plus the eigenvalues of $\frac{1}{\lambda}(I_r\otimes K^{-1})(Z\otimes K)^T SS^T (Z\otimes K)$.

$(Z\otimes K)^T SS^T(Z\otimes K) = (Z^T \otimes K)SS^T(Z \otimes K)$.

$(I_r \otimes K^{-1})(Z^T\otimes K) SS^T (Z\otimes K) = (Z^T \otimes I_n) SS^T (Z\otimes K)$
(using $(I_r\otimes K^{-1})(Z^T\otimes K) = Z^T \otimes (K^{-1}K) = Z^T \otimes I_n$).

And $(Z^T\otimes I_n)SS^T(Z\otimes K) = (Z^T\otimes I_n)\Omega(Z\otimes K)$.

The eigenvalues of this are bounded by $\|Z\|_2^2 \cdot \|K\|_2 \cdot \|\Omega\|_2 = \sigma_{\max}(Z)^2 \cdot \sigma_{\max}(K)$.

So $\kappa(P^{-1}\mathcal{A}) \leq 1 + \sigma_{\max}(Z)^2 \sigma_{\max}(K)/\lambda$.

The regularization $\lambda$ controls the condition number: larger $\lambda$ gives better conditioning but higher approximation error.
\end{proof}

\section{Complexity Analysis}

\begin{theorem}[Complexity of PCG for \eqref{eq:system}]\label{thm:complexity}
Under the assumptions $n, r \ll q \ll N$ (so $n, r < q$):
\begin{enumerate}
\item Each PCG matrix-vector product costs $O(n^2 r + qr)$.
\item Each PCG preconditioner application costs $O(n^2 r + nr^2 + n^3 + r^3)$. The eigendecompositions $O(n^3 + r^3)$ are computed once in precomputation.
\item The number of PCG iterations to achieve tolerance $\varepsilon$ is $O(\sqrt{\kappa} \log(1/\varepsilon))$ where $\kappa = \kappa(P^{-1}\mathcal{A})$.
\item Total cost: $O\!\left((n^3 + r^3 + Mr^2) + \sqrt{\kappa}\log(1/\varepsilon) \cdot (n^2 r + qr)\right)$.
\end{enumerate}

All computations are $O(qr + n^2 r)$ per iteration, avoiding any $O(N)$ computation.
\end{theorem}

\begin{proof}
\textbf{(1) Matrix-vector product:} Derived in Section 4 above. The key is that we never explicitly form the $N$-dimensional vector $(Z\otimes K)v$; instead, we use the sparsity of the observation mask $\Omega = SS^T$ to reduce to $O(q)$ operations.

Specifically:
\begin{itemize}
\item $KV$: $n \times n$ matrix times $n \times r$ matrix = $O(n^2 r)$.
\item Dot products with rows of $Z$: $q$ dot products of $r$-vectors = $O(qr)$.
\item Accumulate $X_{\mathrm{obs}}Z$: $q$ rank-1 updates to $n\times r$ matrix = $O(qr)$.
\item $K(X_{\mathrm{obs}}Z)$: $n\times n$ times $n\times r$ = $O(n^2 r)$.
\end{itemize}

No step requires $O(N)$ or $O(M)$ work (note $M = N/n$, and $n^2 r \ll qr \ll Nr$ since $n < q < N$).

\textbf{(2) Preconditioner application:} $(I_r \otimes K^{-1})v = \mathrm{vec}(K^{-1}V)$ where $V$ is the matrix form of $v$. If $K = V_K \Sigma_K V_K^T$ is precomputed (cost $O(n^3)$ once), then $K^{-1} = V_K \Sigma_K^{-1} V_K^T$ and:
\[
K^{-1}V = V_K (\Sigma_K^{-1} (V_K^T V)).
\]
Cost: $O(n^2 r)$ (three matrix multiplications of $n\times n$ by $n\times r$ matrices).

With the better preconditioner $(Z^TZ) \otimes K^2 + \lambda(I_r\otimes K)$, the eigendecomposition gives: eigenvalues $\lambda_i(Z^TZ) \cdot \sigma_j(K)^2 + \lambda \sigma_j(K)$ for $i \in [r], j \in [n]$. Applying the inverse requires transforming the $nr$-vector by $U_Z \otimes V_K$ (cost $O(n^2 r + nr^2)$), dividing by diagonal entries (cost $O(nr)$), and transforming back (cost $O(n^2 r + nr^2)$).

\textbf{(3) Convergence:} Standard CG convergence bound (Theorem \ref{thm:cg-conv}) gives $O(\sqrt{\kappa}\log(1/\varepsilon))$ iterations for $\varepsilon$-accuracy. With preconditioner, $\kappa = O(1 + \sigma_{\max}(Z)^2\sigma_{\max}(K)/\lambda)$, which is a constant for fixed regularization strength.

\textbf{(4) Total cost:}
\begin{itemize}
\item Precomputation: $K$ eigendecomposition $O(n^3)$, $Z^TZ$ computation $O(Mr^2)$, $Z^TZ$ eigendecomposition $O(r^3)$, MTTKRP $B = TZ$ costs $O(qr)$ (using sparsity of observations).
\item Per-iteration: $O(n^2 r + qr + n^2 r + nr^2) = O(n^2 r + qr + nr^2)$.
\item Total: $O(n^3 + r^3 + Mr^2 + \sqrt{\kappa}\log(1/\varepsilon)(n^2 r + qr))$.
\end{itemize}

No computation of order $O(N)$ appears at any step.

\textbf{Comparison with direct method:} Direct Cholesky on the $nr \times nr$ system costs $O(n^3 r^3)$. PCG costs $O(n^3 + r^3 + Mr^2 + \sqrt{\kappa}\log(1/\varepsilon)(n^2 r + qr))$. Since $n,r \ll q \ll N$, the PCG approach is much faster.

\end{proof}

\section{Why This Works: Summary}

The key ideas enabling the efficient PCG implementation are:

\begin{enumerate}
\item \textbf{Exploiting sparsity of observations:} The matrix $SS^T$ has only $q$ nonzero entries, so matrix-vector products with $(Z\otimes K)^T SS^T (Z\otimes K)$ can be computed in $O(qr)$ time by working with the $q$ observed entries directly, never materializing the full $N$-dimensional vectors.

\item \textbf{Kronecker structure:} The system has Kronecker product structure $(Z\otimes K)^T (\cdots)(Z\otimes K)$, which allows matrix-vector products to be decomposed into smaller operations ($n^2 r$ for the $K$ part, $qr$ for the observation part).

\item \textbf{Good preconditioner:} The preconditioner $P = \lambda(I_r \otimes K)$ or the improved version $(Z^TZ)\otimes K^2 + \lambda(I_r \otimes K)$ can be applied efficiently using precomputed eigendecompositions of $K$ and $Z^TZ$, reducing the condition number and hence the number of CG iterations.

\item \textbf{Avoiding $O(N)$ computation:} At no step do we form the full $N$-dimensional prediction vector or any $N \times \cdot$ matrix. All operations are bounded by $O(qr + n^2 r)$ per iteration.
\end{enumerate}

\section{Verification Protocol}

\textbf{System correctness:} Equation \eqref{eq:system} is the normal equation for the regularized least-squares problem $\min_W \|S^T(Z\otimes K)\mathrm{vec}(W) - S^T\mathrm{vec}(T)\|^2 + \lambda\mathrm{vec}(W)^T(I_r\otimes K)\mathrm{vec}(W)$. This is derived by differentiating the objective and setting to zero.

\textbf{PD of $\mathcal{A}$:} $\mathcal{A} = $ (PSD) + $\lambda$(PD) = PD. Verified.

\textbf{CG convergence theorem:} Standard result \cite[Theorem 6.29]{Trefethen1997}, hypotheses verified.

\textbf{Complexity:} Each step verified above: no $O(N)$ or $O(M)$ (when $M \gg q$) computation.

\textbf{Constraint $n, r < q \ll N$:} Used to ensure $n^2 r + qr \ll N$ (since $qr < q^2 \ll N^2$ and $n^2 r < q^2 r \ll N$).

\begin{thebibliography}{99}

\bibitem{Golub2013}
G.H.\ Golub and C.F.\ Van Loan,
\textit{Matrix Computations}, 4th ed.,
Johns Hopkins University Press, 2013.

\bibitem{Kolda2009}
T.G.\ Kolda and B.W.\ Bader,
\textit{Tensor decompositions and applications},
SIAM Rev.\ \textbf{51} (2009), no.\ 3, 455--500.

\bibitem{Rasmussen2006}
C.E.\ Rasmussen and C.K.I.\ Williams,
\textit{Gaussian Processes for Machine Learning},
MIT Press, 2006.

\bibitem{Trefethen1997}
L.N.\ Trefethen and D.\ Bau III,
\textit{Numerical Linear Algebra},
SIAM, 1997.

\bibitem{SheffieldTheis2008}
N.D.\ Lawrence and J.\ Urtasun,
\textit{Non-linear matrix factorization with Gaussian processes},
Proc.\ ICML, 2009.

\end{thebibliography}

\end{document}
